% ==================================================================
% id:       seq-029
% topic:    סדרות
% subtopic: סדרות מהצורה (1+aₙ)^(1/aₙ)
% solution: full
% ==================================================================
\textbf{חשבו את גבול הסדרות הבאות ונמקו את צעדיכם.}
\begin{enumerate}\bMs{2}
\item $\ds\lim_{n \to \infty} \left( \frac{2n + 2}{2n + 1} \right)^{2n}$
\hint{זהו גבול מהצורה $1^\infty$. כתבו את הבסיס כ-$1+$(ביטוי ששואף לאפס) והשתמשו בגבול המגדיר את $e$.}
\begin{solbox}
בגלל שזהו מצב מהצורה
$a_n^{b_n}$
כאשר
$a_n\goes 1$
ו-$b_n\goes\infty$,
אין לנו כלל שאומר לאן תתכנס הסדרה
$a_n^{b_n}$,
ונצטרך לעבוד קצת.
נציג
\begin{align*}
    \left( \frac{2n + 2}{2n + 1} \right)^{2n}&=\
    \left( \frac{2n + 1+1}{2n + 1} \right)^{2n}=\
    \left(1+ \frac{1}{2n + 1} \right)^{2n}\\
    &=\ \left(1+ \frac{1}{2n + 1} \right)^{2n{\color{red}+1-1}}\\
    &=\ \frac{1}{1+ \frac{1}{2n + 1}}\cdot\left(1+ \frac{1}{2n + 1} \right)^{2n+1}
\end{align*}
נזכור את העובדה שנלמדה בכיתה שלפיה
$$,\lim_{n\to\infty}(1+x_n)^{\frac{1}{x_n}}=e$$
כאשר
$(x_n)_{n=1}^\infty$
סדרה המקיימת
$\displaystyle\lim_{n\to\infty}x_n=0$.
במקרה הנוכחי
$x_n=\frac{1}{2n+1}$.
לכן הגבול של המונה הוא
$e$,
מכאן,
לפי אריתמטיקה של גבולות (כי הגבול של המכנה הוא
$1$), נסיק כי הגבול המבוקש הוא
$e$.
\end{solbox}
%%%%%
\item $\ds\lim_{n \to \infty} \lr{\frac{3n^2+4n-2}{3n^2-6n-5}}^{\frac{n^2+3}{2n+5}}$
\hint{זהו גבול מהצורה $1^\infty$. כתבו את הבסיס כ-$1+$(ביטוי ששואף לאפס) וסדרו את החזקה כך שתופיע בה ההופכית של אותו ביטוי.}
\begin{solbox}
נציג
\begin{align*}
    \left(\frac{3n^2+4n-2}{3n^2-6n-5}\right)^{\frac{n^2+3}{2n+5}}&=\ \left(\frac{3n^2{\color{red}-6n-5}+4n-2{\color{red}+6n+5}}{3n^2-6n-5}\right)^{\frac{n^2+3}{2n+5}}\\
    &=\ \left(1+\frac{10n+3}{3n^2-6n-5}\right)^{\frac{n^2+3}{2n+5}}\\
    &=\ \left(1+\frac{10n+3}{3n^2-6n-5}\right)^{{\color{red}\frac{3n^2-6n-5}{10n+3}}\cdot\frac{n^2+3}{2n+5}\cdot{\color{red}\frac{10n+3}{3n^2-6-5n}}}.
\end{align*}
ראינו בכיתה כי
$$,\displaystyle\lim_{n\to\infty}\left(1+\frac{10n+3}{3n^2-6n-5}\right)^{\frac{3n^2=6n-5}{10n+3}}=e$$
בנוסף, לא קשה להשתכנע כי
$$,\displaystyle\lim_{n\to\infty}\frac{n^2+3}{2n+5}\cdot\frac{10n+3}{3n^2-6-5n}=\frac{5}{3}$$
לכן מאריתמטיקה של חזקות, הגבול המבוקש הוא
$e^{\frac{5}{3}}$.
\end{solbox}
%%%%%
\item $\ds\lim_{n \to \infty} \lr{1-\frac{1}{n^2}}^{n}$
\hint{זהו גבול מהצורה $1^\infty$. כתבו את הבסיס כ-$1+$(ביטוי ששואף לאפס) וסדרו את החזקה כך שתופיע בה ההופכית של אותו ביטוי, כמו בסעיפים הקודמים.}
\begin{solbox}
כמו בסעיפים הקודמים, זהו מצב של
$"1^{\infty}"$.
כאן $x_n=-\frac{1}{n^2}\goes0$, ולכן נסדר את החזקה כך שתופיע בה ההופכית $\frac{1}{x_n}=-n^2$:
$$.\lr{1-\frac{1}{n^2}}^{n}=\lr{\lr{1-\frac{1}{n^2}}^{{\color{red}-n^2}}}^{{\color{red}-\frac{1}{n^2}}\cdot n}=\lr{\lr{1-\frac{1}{n^2}}^{-n^2}}^{-\frac{1}{n}}$$
לפי העובדה שנלמדה בכיתה, $\lr{1+x_n}^{\frac{1}{x_n}}\goes e$, כלומר הבסיס שואף ל-$e$. המעריך $-\frac1n$ שואף ל-$0$.
לכן, לפי אריתמטיקה של חזקות, הגבול המבוקש הוא
$e^0=1$.
\end{solbox}
%%%%%
\item $\ds\lim_{n \to \infty} n\lr{\ln(n+3)-\ln n}$
\hint{השתמשו בחוקי הלוגריתמים כדי לכתוב את הסדרה כ-$\ln$ של ביטוי מהצורה $1^\infty$, כמו בסעיפים הקודמים.}
\begin{solbox}
לפי חוקי הלוגריתמים,
$$.n\lr{\ln(n+3)-\ln n}=n\ln\lr{\frac{n+3}{n}}=\ln\lr{1+\frac3n}^{n}$$
בתוך ה-$\ln$ יש מצב של
$"1^{\infty}"$,
וכמו בסעיפים הקודמים
$$,\lr{1+\frac3n}^{n}=\lr{\lr{1+\frac3n}^{\frac n3}}^{3}\goes e^3$$
לפי העובדה $\lr{1+x_n}^{\frac{1}{x_n}}\goes e$ עם $x_n=\frac3n$, ולפי אריתמטיקה של חזקות.
כדי לעבור לגבול של ה-$\ln$ בלי להשתמש ברציפות, נחסום. לכל $x>0$ מתקיים $\frac{x}{1+x}\le\ln(1+x)\le x$. עם $x=\frac3n$ נקבל
$$.\frac{3}{1+\frac3n}=n\cdot\frac{\frac3n}{1+\frac3n}\le n\ln\lr{1+\frac3n}\le n\cdot\frac3n=3$$
אגף שמאל שואף ל-$3$ לפי אריתמטיקה של גבולות, ולכן לפי כלל הסנדוויץ' הגבול המבוקש הוא
$3$
(בהתאם לכך ש-$\ln e^3=3$).
\end{solbox}
\eMs
\end{enumerate}
